\section{Conclusiones}

Esta revisión consolida a las nanopartículas de plata (AgNPs) como un sistema
físico fundamental en la intersección de la nanofotónica y la ciencia de
biomateriales. La fenomenología de las AgNPs emerge de la resonancia de
plasmón de superficie localizado, una excitación electrónica colectiva
estrictamente parametrizada por la geometría de la frontera y la permitividad
dieléctrica del entorno. El confinamiento electromagnético resultante amplifica
de manera masiva la susceptibilidad óptica no lineal. La cuantificación
experimental de esta dinámica mediante la técnica de escaneo Z corrobora su
viabilidad para la limitación óptica pasiva y la conmutación fotónica
ultrarrápida. De manera concurrente, la síntesis biogénica otorga una ruta
termodinámica que mitiga la citotoxicidad de los precursores convencionales,
estabilizando los nanomateriales mediante biocoronas que modulan sus
interacciones biológicas y optofísicas.

No obstante, la explotación tecnológica de estos sistemas requiere solucionar
limitaciones físicas y experimentales críticas. Primero, la estocasticidad
bioquímica de los métodos verdes demanda modelos cinéticos rigurosos que
garanticen el control determinista sobre la polidispersidad morfológica, la
variable macroscópica que gobierna la homogeneidad de la respuesta óptica.
Segundo, la electrodinámica del acoplamiento plasmónico en arquitecturas
heterogéneas, tales como interfaces con semiconductores o materiales
bidimensionales, exige formulaciones teóricas que superen la aproximación
cuasiestática de dispersores aislados. Finalmente, la integración de estas
nanoestructuras en dispositivos optoelectrónicos macroscópicos y biosensores
está condicionada a la medición precisa de su estabilidad termodinámica, sus
umbrales de fotodegradación bajo campos láser intensos y su lixiviación iónica
en diferentes matrices. La resolución de estos problemas de frontera
establecerá el rigor necesario para el diseño teórico y experimental de
biocompuestos fotónicos de nueva generación.
